\chapter{Unitary Dynamics and Schr\"odinger's Equation}
\label{ch:dynamics}

\section{An Explicit New Assumption: Continuous Reversible Evolution}

Parts I--II established that coherent, reversible transformations between
\emph{discrete} contexts are represented by unitary lifts
(Chapters~\ref{ch:composition}--\ref{ch:real-mub-obstruction}). Nothing
proved so far concerns physical \emph{time}. Connecting the two requires a
genuinely new assumption:

\begin{assumption}[Continuous reversible time evolution]
\label{ass:continuous-evolution}
A closed system's evolution over time is represented by a strongly
continuous one-parameter family of unitary operators $U(t)$ satisfying
\begin{equation}
U(t+s) = U(t)U(s), \qquad U(0) = I.
\end{equation}
\end{assumption}

This is not derived from the discrete admissibility framework of Parts I--II
— it imports continuity in time as new physical content, using "unitary"
only because Parts I--II already established unitary operators as the
coherent lifts of admissible transformations. The group law itself is the
same composition law studied throughout the book (Chapter
\ref{ch:context-graphs}'s exact closure, Chapter~\ref{ch:composition}'s
composition rule), now imposed on a continuous parameter rather than a
discrete sequence of contexts.

\section{The Generator Theorem}

In finite dimension, Assumption~\ref{ass:continuous-evolution} has an
elementary consequence — no need for the general (infinite-dimensional)
machinery of Stone's theorem.

\begin{theorem}[Finite-dimensional generator theorem]
\label{thm:generator}
\index{Finite-dimensional generator theorem (theorem)}
Given Assumption~\ref{ass:continuous-evolution} with $U(t)$ differentiable at
$t=0$, there is a Hermitian operator $H$ and real constant $\hbar$ such that
\begin{equation}
U(t) = e^{-itH/\hbar}.
\end{equation}
\end{theorem}

\begin{proof}
Let $K := \frac{d}{dt}U(t)\big|_{t=0}$. Differentiating $U(t+s)=U(t)U(s)$
with respect to $s$ at $s=0$ gives $\frac{d}{dt}U(t) = U(t)K$, a linear ODE
with solution $U(t) = e^{Kt}$ (using $U(0)=I$). Differentiating
$U(t)^\dagger U(t) = I$ at $t=0$ gives $K^\dagger + K = 0$, i.e.\ $K$ is
anti-Hermitian, so $K = -iH/\hbar$ for some Hermitian $H$ and real scale
$\hbar$ (any decomposition of an anti-Hermitian operator as $-i$ times a
Hermitian one, with the scale factor absorbed into $H$ or left explicit as
$\hbar$).
\end{proof}

\begin{remark}
This is genuinely elementary — a finite-dimensional linear ODE argument, not
an import of Stone's theorem (which handles the harder infinite-dimensional
case this book's finite-$d$ scope does not need). The differentiability
assumed at $t=0$ is itself part of Assumption~\ref{ass:continuous-evolution}'s
content; strong continuity alone, without differentiability, requires the
fuller argument this book does not attempt.
\end{remark}

\section{Schr\"odinger's Equation}

\begin{corollary}[Schr\"odinger's equation]
\label{cor:schrodinger}
\index{Schr\"odinger's equation (corollary)}
Writing $|\psi(t)\rangle := U(t)|\psi(0)\rangle$,
\begin{equation}
i\hbar\,\frac{d}{dt}|\psi(t)\rangle = H|\psi(t)\rangle.
\end{equation}
\end{corollary}

\begin{proof}
$\frac{d}{dt}|\psi(t)\rangle = \frac{d}{dt}U(t)|\psi(0)\rangle = -\tfrac{i}{
\hbar}H\,U(t)|\psi(0)\rangle = -\tfrac{i}{\hbar}H|\psi(t)\rangle$, using
Theorem~\ref{thm:generator} and $\frac{d}{dt}e^{-itH/\hbar} = -\tfrac{i}{
\hbar}H\,e^{-itH/\hbar}$ (differentiating the matrix exponential). Multiplying
by $i\hbar$ gives the stated equation.
\end{proof}

\begin{remark}
Schr\"odinger's equation is, on this route, the infinitesimal form of the
composition law of Assumption~\ref{ass:continuous-evolution} — not an
independent postulate stacked on top of the state-space structure recovered
in Chapters~\ref{ch:born-rule}--\ref{ch:observables}.
\end{remark}

\section{What \texorpdfstring{$H$}{H} Is Not (Yet)}

Theorem~\ref{thm:generator} produces \emph{a} Hermitian generator and \emph{a}
real scale $\hbar$. It does not say that $H$ is energy, nor that $\hbar$ is
Planck's constant. Those identifications are additional physical input —
typically made via correspondence with classical Hamiltonian mechanics (the
generator of time translation is identified with energy through Noether's
theorem or the classical Hamiltonian's role as generator of time evolution in
the classical limit) — and are not established by anything in this chapter.
The generator theorem gives the \emph{form} $i\hbar\,d_t|\psi\rangle =
H|\psi\rangle$; it does not by itself explain why the object called $H$
deserves to be called energy.

\begin{ledgerentry}[End of Chapter~\ref{ch:dynamics}]
\textbf{Established:} given Assumption~\ref{ass:continuous-evolution} (stated
explicitly as new, not derived), the finite-dimensional generator theorem
(Theorem~\ref{thm:generator}) and Schr\"odinger's equation as its
infinitesimal form (Corollary~\ref{cor:schrodinger}), via an elementary
matrix-ODE argument rather than the general Stone's theorem.\\
\textbf{Not established:} that physical time evolution actually satisfies
Assumption~\ref{ass:continuous-evolution} — this is imported, not derived
from Parts I--II's discrete admissibility framework; that the generator $H$
is energy or that $\hbar$ is Planck's constant, both of which require
additional physical identification not attempted here; anything about
open-system or non-unitary dynamics, which would require the channel/Kraus
formalism this book has not developed.
\end{ledgerentry}
